% Round 003 restricted-model upper bounds; not bounds for unrestricted C_n.
\subsection{任意集合的同半径部分 halo}
\label{model:partialhalo}
固定 $n\ge2$，记
\[
B_n=2^n,\qquad p=1+\frac1n,\qquad a_n=\frac{n}{2(n+1)},\qquad A_n(t)=1+\frac{B_n^2t}{1-B_n t}\quad(0<t<1/B_n).
\]
令 $Q(x,r)=x+[-r,r]^n$、$H_s(E)=\{x:M_{c,n}\ind_E(x)>1-s\}$。
\begin{proposition}[任意 $E$ 的同半径部分 halo]
设 $0<|E|<\infty$，$F\subset E^c$ 可测，且对某个共同半径 $r>0$， 每个 $x\in F$ 满足
\begin{equation}\label{eq:partial-halo-witness}
|Q(x,r)\setminus E|<t(2r)^n.
\end{equation}
置 $z=|F|/|E|$。则
\begin{equation}\label{eq:partial-halo-bound}
\frac{z}{A_n(t)}+a_n(1-B_n t)^p z^p\le B_n t.
\end{equation}
特别地，令 $\kappa_n=a_nB_n^p$，当 $t\downarrow0$ 时
\begin{equation}\label{eq:partial-halo-endpoint}
z\le B_n t-\kappa_nt^p+O_n(t^{1+2/n}+t^2).
\end{equation}
这是同半径部分 halo 的界，并非无约束 $C_n(1-t)$ 的上界。
\end{proposition}
\begin{proof}
若 $|F|=0$ 则平凡。又 $F\subset H_t(E)$，故基线上界保证 $|F|<\infty$。令 $J_r=\ind_{(-r/2,r/2)^n}$，
\[
D=F+(-r/2,r/2)^n.
\]
若 $y\in D$，存在 $x\in F$ 使 $\|x-y\|_\infty<r/2$，故 $Q(y,r/2)\subset Q(x,r)$。由 \eqref{eq:partial-halo-witness}，
\[
|Q(y,r/2)\setminus E|<t(2r)^n=B_n t|Q(y,r/2)|,
\]
所以 $D\subset H_{B_n t}(E)$。基线上界 $|H_s(E)|\le U_n(s)|E|$ 给
\begin{equation}\label{eq:partial-halo-D}
0<|D|\le U_n(B_n t)|E|=A_n(t)|E|.
\end{equation}
这里 $D$ 包含一个非空开方盒，故其测度严格为正。 记 $S=\ind_F*J_r$。其支撑包含于 $D$，因为 $S(u)>0$ 只可能在 存在 $x\in F$ 满足 $u-x\in(-r/2,r/2)^n$ 时发生。又 $\int S=|F|r^n$，从而 Tonelli 与 Cauchy--Schwarz 给
\begin{equation}\label{eq:partial-halo-triangle}
\mathcal E_T:=\iint_{F\times F}T_r(x-y)\,dx\,dy =r^{-n}\int_D S^2 \ge r^n\frac{|F|^2}{A_n(t)|E|},
\end{equation}
其中 $T_r(h)=r^{-n}(J_r*J_r)(h)=\prod_j(1-|h_j|/r)_+$。 将 $\mathbb R^n$ 分为边长 $r$ 的半开方格，只保留与 $F$ 正测交叠的活动格。 若 $P$ 是活动格，则取 $x\in P\cap F$，有 $P\subset Q(x,r)$，于是 $|E\cap P|>(1-B_n t)r^n$。活动格两两不交且 $|E|<\infty$，故其数量 $N$ 有限，并满足
\begin{equation}\label{eq:partial-halo-grid}
Nr^n\le\frac{|E|}{1-B_n t}.
\end{equation}
令 $K_r=\ind_{[-r,r]^n}$。在 $[-r,r]^n$ 上，置 $u_j=|h_j|/r$，则 $K_r(h)-T_r(h)=1-\prod_j(1-u_j)\ge\max_j u_j=\|h\|_\infty/r$。 若 $B\subset\mathbb R^n$、$|B|=b$，层蛋糕公式及 $|B\cap Q(x,u)|\le(2u)^n$ 给
\begin{equation}\label{eq:distance-moment}
\int_B\|x-y\|_\infty\,dy \ge\int_0^{b^{1/n}/2}[b-(2u)^n]\,du=a_nb^p, \qquad \iint_{B\times B}\|x-y\|_\infty\,dx\,dy\ge a_nb^{1+p}.
\end{equation}
令 $b_j=|F\cap P_j|$。同格点对、 \eqref{eq:partial-halo-grid}、\eqref{eq:distance-moment} 与 Jensen 不等式给
\[
\mathcal E_{K-T}:=\iint_{F\times F}(K_r-T_r)(x-y)\,dx\,dy \ge a_nr^{-1}\sum_jb_j^{1+p} \ge a_nr^n(1-B_n t)^p\frac{|F|^{1+p}}{|E|^p}.
\]
另一方面，由 $F\subset E^c$ 及 \eqref{eq:partial-halo-witness}，
\[
\mathcal E_K:=\iint_{F\times F}K_r(x-y)\,dx\,dy =\int_F|F\cap Q(x,r)|\,dx\le B_n t\,r^n|F|.
\]
合并 $\mathcal E_K=\mathcal E_T+\mathcal E_{K-T}$ 与 \eqref{eq:partial-halo-triangle}，再除以 $r^n|F|$，得 \eqref{eq:partial-halo-bound}。 左端关于 $z$ 严格递增。设 $z_*(t)$ 为等号的唯一非负根。利用 $A_n(t)=1+O_n(t)$、$(1-B_n t)^p=1+O_n(t)$，代入根方程先得 $z_*(t)=B_n t+O_n(t^p)$，再代回一次得 $z_*(t)=B_n t-a_nB_n^pt^p+O_n(t^{1+2/n}+t^2)$，即 \eqref{eq:partial-halo-endpoint}。\qedhere
\end{proof}
\begin{remark}[近同半径带]
若各 $x\in F$ 的见证半径 $r_x\in[r,(1+\eta)r]$，且 $|Q(x,r_x)\setminus E|<t(2r_x)^n$，则 $Q(x,r)\subset Q(x,r_x)$， 故命题以 $s=t(1+\eta)^n$ 替代 $t$ 仍成立，只需 $B_n s<1$。 当 $\eta=o(t^{1/n})$ 时，它给 $z\le B_n t-\kappa_nt^p+o_n(t^p)$；此处一般只有小 $o$ 余项。
\end{remark}
\subsection{Lipschitz 见证半径下的全局可变尺度界}
\label{model:lipschitz}
令 $0<|E|<\infty$，$F\subset E^c$ 可测。假设存在可测选择 $r:F\to(0,\infty)$，使每个 $x\in F$ 满足
\begin{equation}\label{eq:lipschitz-witness}
|Q(x,r(x))\setminus E|<t(2r(x))^n, \qquad |r(x)-r(y)|\le L\|x-y\|_\infty\quad(x,y\in F).
\end{equation}
记 $z=|F|/|E|$、$\lambda=(1+L/2)^{-1}$、$C=2(1+L)$。
\begin{theorem}[缓变见证半径的部分 halo 界]
若 $(2C)^nt<1$，则
\begin{equation}\label{eq:lipschitz-bound}
\frac{\lambda^n z}{A_n(t)}+ a_nC^{-(n+1)}[1-(2C)^nt]^p z^p\le B_n t.
\end{equation}
特别地，若 $L=L_t=o(t^{1/n})$，则
\begin{equation}\label{eq:lipschitz-endpoint}
z\le B_n t-a_nt^p+o_n(t^p).
\end{equation}
此为具有 Lipschitz 见证选择的受限类结论，不是无约束 $C_n(1-t)$ 的上界。
\end{theorem}
\begin{proof}
$F$ 零测时平凡。以下设 $|F|>0$；由 $F\subset H_t(E)$ 及基线上界，$|F|<\infty$。对 $x\in F$ 定义开方盒
\[
B_x=x+(-\lambda r(x)/2,\lambda r(x)/2)^n,\quad \phi_x=\frac{\ind_{B_x}}{(\lambda r(x))^n},\quad G(x,y)=\int\phi_x(u)\phi_y(u)\,du.
\]
Gram 表示给出 $G\ge0$。记 $d=\|x-y\|_\infty$、$r=r(x)$、$s=r(y)$。 下证对有序的每一对点均有
\begin{equation}\label{eq:gram-pair}
G(x,y)\le\lambda^{-n}r(x)^{-n} \left(1-\frac{d}{r(x)}\right)_+.
\end{equation}
若 $d=0$，由 Lipschitz 条件 $r=s$ 且两盒重合，\eqref{eq:gram-pair} 取等。设 $d>0$，并先假设 $r\ge s$。Lipschitz 条件给
\[
\frac{\lambda|r-s|}{2}\le\frac{\lambda Ld}{2}<d,
\]
故两盒互不包含。若交集有正测度，必有 $d<\lambda(r+s)/2\le\lambda(s+Ld/2)$；因 $1-\lambda L/2=\lambda$，得 $d<s$。 在达到中心距离的坐标，交叠长度至多 $\lambda(r+s)/2-d$；其余坐标至多 $\lambda s$。因此
\[
G(x,y)\le\lambda^{-n}r^{-n} \left[1+\frac{r-s}{2s}-\frac{d}{\lambda s}\right] \le\lambda^{-n}r^{-n}(1-d/s) \le\lambda^{-n}r^{-n}(1-d/r),
\]
其中第二步用 $r-s\le Ld$ 及 $\lambda^{-1}-L/2=1$，末步用 $s\le r$。 若交集零测则结论平凡。中间界还满足 $r^{-n}(1-d/s)\le s^{-n}(1-d/s)$，故也控制小半径一侧的有序行。交换 $x,y$ 可处理 $r<s$，并同时说明 $G(x,y)>0$ 必有 $d<\min(r,s)$；故当 $d\ge r(x)$ 时 $G=0$， \eqref{eq:gram-pair} 对全体有序点成立。 置
\[
D(x,y)=\lambda^{-n}r(x)^{-n}\ind_{\{d\le r(x)\}}-G(x,y)\ge0.
\]
由 \eqref{eq:gram-pair}，当 $d\le r(x)$ 时
\begin{equation}\label{eq:gram-deficit}
D(x,y)\ge\lambda^{-n}r(x)^{-(n+1)}d.
\end{equation}
原见证及 $F\subset E^c$ 给
\begin{equation}\label{eq:lipschitz-I}
I:=\int_F r(x)^{-n}|F\cap Q(x,r(x))|\,dx\le B_n t|F|.
\end{equation}
令 $S(u)=\int_F\phi_x(u)\,dx$、$D_0=\bigcup_{x\in F}B_x$。 $S$ 支撑在 $D_0$ 内、$\int S=|F|$；而若 $u\in B_x$，则 $Q(u,r(x)/2)\subset Q(x,r(x))$，故 $D_0\subset H_{B_n t}(E)$。 基线上界给 $0<|D_0|\le A_n(t)|E|$，从而
\begin{equation}\label{eq:lipschitz-gram}
\iint_{F\times F}G(x,y)\,dx\,dy=\int_{D_0}S^2 \ge\frac{|F|^2}{A_n(t)|E|}.
\end{equation}
取一个 dyadic 网格。称半开 dyadic 格 $P$ 为好格，若 $\ell(P)\le\inf_{x\in P\cap F}r(x)$，空交集的下确界视为 $+\infty$。 每个 $x\in F$ 的充分小格均为好格：当 $y\in P\cap F$ 时 $r(y)\ge r(x)-L\ell(P)$。充分大的祖格不是好格，因为其中含 $x$ 且 $\ell(P)>r(x)$。好格向下遗传，故取所有与 $F$ 正测交叠的极大好格， 它们两两不交并覆盖 $F$，至多差零测集。对极大好格 $P$，令 $\ell=\ell(P)$； 父格不好，故存在 $y\in P^+\cap F$ 使 $r(y)<2\ell$。对 $x\in P\cap F$，
\[
\ell\le r(x)\le r(y)+2L\ell<C\ell.
\]
取 $P\cap F$ 中一个点的见证，$P\subset Q(x,r(x))$，因而 $|E\cap P|>[1-(2C)^nt]|P|$。若 $V_* =\sum_P|P|$ 为活动格总体积， 即使活动格可数无穷，也有
\begin{equation}\label{eq:lipschitz-active-grid}
V_*\le\frac{|E|}{1-(2C)^nt}.
\end{equation}
对每个格 $P$ 记 $b_P=|F\cap P|$。同格点距满足 $d\le\ell\le r(x)$，故由 \eqref{eq:gram-deficit}、\eqref{eq:distance-moment} 和 $r(x)<C\ell$， 再对格体积作带权 Jensen，
\[
 \begin{aligned}
 \iint_{F\times F}D(x,y)\,dx\,dy
 &\ge\lambda^{-n}a_nC^{-(n+1)}
       \sum_P\frac{b_P^{1+p}}{|P|^p}\\
 &\ge\lambda^{-n}a_nC^{-(n+1)}
       \frac{|F|^{1+p}}{V_*^p}\\
 &\ge\lambda^{-n}a_nC^{-(n+1)}
       [1-(2C)^nt]^p\frac{|F|^{1+p}}{|E|^p}.
 \end{aligned}
\]
按定义 $\lambda^{-n}I=\iint G+\iint D$。合并 \eqref{eq:lipschitz-I}、\eqref{eq:lipschitz-gram} 与上式，乘以 $\lambda^n/|F|$ 即得 \eqref{eq:lipschitz-bound}。被积量非负，故可用 Tonelli；\eqref{eq:lipschitz-I} 也保证自能有限。 最后令 $u=t^{1/n}$。若 $L=o(u)$，则 $\lambda^n/A_n(t)=1+o(u)$、$C=2+o(1)$、 $[1-(2C)^nt]^p=1+O_n(t)$。令 $z_*(t)$ 为 \eqref{eq:lipschitz-bound} 左端等于 $B_n t$ 的唯一非负根。记 $\alpha_t=\lambda^n/A_n(t)=1+o(u)$、 $\beta_t=a_nC^{-(n+1)}[1-(2C)^nt]^p=a_n2^{-(n+1)}+o(1)$，代入根方程得
\[
z_*(t)=\frac{B_n t}{\alpha_t}-\frac{\beta_t}{\alpha_t} \left(\frac{B_n t}{\alpha_t}\right)^p+O_n(t^{2p-1}).
\]
因 $2p-1=1+2/n>p$，且 $\alpha_t=1+o(t^{1/n})$，
\[
z\le z_*(t)=B_n t-a_n2^{-(n+1)}B_n^pt^p+o_n(t^p) =B_n t-a_nt^p+o_n(t^p),
\]
因为 $a_n2^{-(n+1)}B_n^p=a_n$；且 $t^2=o(t^p)$ 对 $n\ge2$。
\end{proof}
